
\section*{Appendix A — Cleaning Procedure}

\begin{center}
\small
\begin{tabular}{llll}
\toprule
Step & Script & Input & Output \\
\midrule
A & \texttt{01 FH Cleaning.sas}             & FH raw           & \texttt{fh\_latest.csv}, \texttt{unique\_ids.csv} \\
B & \texttt{02 MMO Cleaning.sas}            & MMO raw          & \texttt{mmo\_latest.csv} \\
C & \texttt{01 build shortage list.R}       & PDF \textit{arrêtés}, DARES xls & \texttt{shortage\_list.csv} \\
D & \texttt{03 Shortage List + FAP PCS Match.R} & \texttt{shortage\_list.csv}, PCS-ESE ods & \texttt{fap\_to\_pcs.csv} \\
E & \texttt{04 Parquet + Final Cleaning.R}  & all CSVs         & \texttt{mmo\_latest\_wage.parquet} \\
F & \texttt{01 DiD.R}                       & all parquets     & analysis dataset \\
\bottomrule
\end{tabular}
\end{center}

\subsection*{A. FH Cleaning \textnormal{\small(\texttt{01 FH Cleaning.sas})}}

\begin{itemize}
  \item \textbf{Nationality restriction.} Retain only spells where \texttt{NATION}
  belongs to the following codes, corresponding to our four matched pairs plus French nationals:
  \begin{center}
  \small
  \begin{tabular}{lll}
  \toprule
  Code & Country & Group \\
  \midrule
  01 & France & French (reference) \\
  24 & Serbia / Bosnia-Herz. & Non-EEA \\
  26 & Kosovo / Montenegro / N. Macedonia & Non-EEA \\
  29 & Turkey & Non-EEA \\
  27 & Croatia & EEA \\
  96 & Slovenia & EEA \\
  97 & Bulgaria & EEA \\
  98 & Romania & EEA \\
  19 & Greece & EEA \\
  02, 91--95 & Baltic and Central EU & EEA \\
  \bottomrule
  \end{tabular}
  \end{center}
  All other nationalities are discarded to limit intermediate file size.

  \item \textbf{Temporal restriction.} Keep spells with \texttt{datins} $\geq$
  1 January 2019 to ensure at least two pre-treatment years before the April 2021 reform.

  \item \textbf{Education recoding.} \texttt{NIVFOR} alphanumeric codes
  (\texttt{AFS}, \texttt{CP4}, \ldots, \texttt{NV1}) are mapped to an ordered
  integer scale 0--9 (\texttt{formation}) for use as a continuous control.

  \item \textbf{One row per worker.} FH records unemployment spells; multiple
  rows may exist per \texttt{id\_force}.
  \begin{itemize}
    \item Workers with a single spell: kept as-is.
    \item Workers with multiple spells: the row with the latest \texttt{datins}
    is retained, on the grounds that the most recent registration reflects
    current characteristics.
  \end{itemize}

  \item \textbf{ID export.} Distinct \texttt{id\_force} values are exported to
  \texttt{unique\_ids.csv} and used as a hash reference in Step B.
\end{itemize}

\subsection*{B. MMO Cleaning \textnormal{\small(\texttt{02 MMO Cleaning.sas})}}

\begin{itemize}
  \item \textbf{ID filter.} A hash join on \texttt{unique\_ids} restricts MMO
  to workers present in FH. This is necessary to attach nationality and
  controls from FH; workers not registered with France Travail are excluded.

  \item \textbf{Contract filters} (applied to each annual file 2019--2025):
  \begin{itemize}
    \item Drop rows with missing \texttt{L\_Contrat\_SQN}: these cannot be
    tracked across vintages.
    \item Drop contracts starting before 2019.
    \item Drop contracts shorter than 31 days: very short spells do not
    reflect stable employment and would distort monthly wage comparisons.
  \end{itemize}

  \item \textbf{Deduplication across annual files.} The same contract can
  appear in multiple annual MMO files (once per year until its end date),
  with later files potentially containing corrections. A contract-level
  index is built across all years; for each \texttt{L\_Contrat\_SQN},
  only the most recent vintage is kept.

  \item \textbf{Policy timing dummies.} Binary indicators are created for
  contract start after each reform date: \texttt{post\_apr2021} (1 Apr 2021),
  \texttt{post\_jan2024} (29 Jan 2024), \texttt{post\_mar2024} (1 Mar 2024),
  \texttt{post\_apr2025} (1 Apr 2025), \texttt{post\_may2025} (21 May 2025).
  Symmetric two-year DSN recency windows are also created for robustness checks.
\end{itemize}

\subsection*{C--D. Shortage List Construction and FAP--PCS Matching
\textnormal{\small(\texttt{01 build shortage list.R},
\texttt{03 Shortage List + FAP PCS Match.R})}}

\begin{itemize}
  \item The four \textit{arrêtés} (2008, 2021, 2024, 2025) are parsed from
  PDF. Each FAP code is mapped to PCS-2003 codes via the DARES
  correspondence table.
  \item The 2024 amendment is applied as a set of targeted removals and
  additions to the 2021 list. Delisting and relisting flags are constructed.
  \item The DARES correspondence table \texttt{Table\_passage\_FAP\_PCS-ESE.ods}
  is cleaned and exported as \texttt{fap\_to\_pcs.csv}, which maps each FAP
  code to its PCS-ESE codes. This file drives the FAP--PCS join in Step F.
  A verification step checks that all PCS codes in \texttt{shortage\_list.csv}
  appear in this crosswalk; missing codes are flagged in-session.
  \item \textbf{Limitation.} The version of \texttt{Table\_passage\_FAP\_PCS-ESE.ods}
  available on CASD is incomplete: several FAP codes present in the
  \textit{arrêtés} have no entry in the table. Contracts in those occupations
  receive no FAP match in Step F and are silently assigned to the control group,
  causing an undercount of treated contracts.
  \item The FAP--PCS crosswalk and the treatment status assignment are
  \textbf{not merged at this stage}; they are joined to the contracts
  in Step F.
\end{itemize}

\subsection*{E. Wage Harmonisation and Geographic Scope
\textnormal{\small(\texttt{04 Parquet + Final Cleaning.R})}}

\begin{itemize}
  \item \textbf{Format conversion.} All CSVs are converted to Parquet.
  Subsequent operations run via DuckDB without loading files into memory.

  \item \textbf{Wage harmonisation.} A standardised \texttt{monthly\_wage}
  is constructed from \texttt{salaire\_base}, \texttt{ModeExercice}, and
  \texttt{salaire\_base\_mois\_complet}:
  \begin{itemize}
    \item \textit{Part-time}: \texttt{ModeExercice} $\in \{20,30,40,41,42\}$,
    or \texttt{ModeExercice = 99} with \texttt{salaire\_base} $< 1{,}498$
    (2019 SMIC monthly gross).
    \item If part-time and \texttt{salaire\_base\_mois\_complet = 1}:
    \texttt{monthly\_wage = salaire\_base} (already full-month equivalent).
    \item If full-time and \texttt{salaire\_base\_mois\_complet = 0}:
    \texttt{monthly\_wage = salaire\_base} $\times 30.4375\,/\,\text{contract days}$.
    \item All other combinations: set to missing and excluded from regressions.
  \end{itemize}

  \item \textbf{Geographic restriction.} Contracts are assigned to the
  13 metropolitan regions via a postal-code prefix crosswalk. Contracts
  with \texttt{CP} mapping to \textit{Outre-Mer} (prefixes 97--98) or
  abroad are dropped: the shortage lists apply only to metropolitan France.
\end{itemize}

\subsection*{F. Final Merge and Treatment Assignment
\textnormal{\small(\texttt{01 DiD.R})}}

\begin{itemize}
  \item MMO and FH are joined on \texttt{id\_force} (INNER JOIN).
  \item The FAP--PCS crosswalk is joined on \texttt{PcsEse = pcs\_code}
  (LEFT JOIN): contracts whose PCS code is absent from the crosswalk receive
  no FAP match and are assigned to the control group via \texttt{COALESCE}.
  \item The shortage list is joined on \texttt{pcs\_code} and
  \texttt{region\_contract = region} (LEFT JOIN): treatment indicators
  are \texttt{COALESCE}d to 0 for unmatched contracts, which form
  the control group.
  \item Age (\texttt{contract\_year} $-$ birth year), a male dummy,
  a diploma dummy, and experience in months are derived at this stage.
\end{itemize}